\begingroup
\setlength{\parindent}{0pt}
\setlength{\parskip}{0.58\baselineskip}
\setlength{\abovedisplayskip}{9pt plus 2pt minus 2pt}
\setlength{\belowdisplayskip}{9pt plus 2pt minus 2pt}
\setlength{\jot}{7pt}

\section*{How the Proof Works}

We are asking whether a viscous fluid that starts smooth can ever stop being
smooth.  This proof answers no.  Start with a divergence-free velocity \(u_0\)
on \(\mathbb R^3\) that is smooth and rapidly decaying with all its derivatives;
this is the datum class \(\mathcal S_\sigma\).  Keep the viscosity \(\nu>0\)
fixed.  The difficulty is that viscosity is not the
only part of the equation changing the fluid.  The fluid also carries and
deforms its own velocity:
\[
\underbrace{\partial_tu}_{\text{local acceleration}}
+\underbrace{(u\!\cdot\nabla)u}_{\text{nonlinear transport}}
=-\underbrace{\nabla p}_{\text{pressure}}
+\underbrace{\nu\Delta u}_{\text{viscous diffusion}},
\qquad \nabla\!\cdot u=0.
\]
The nonlinear term can sharpen the flow while viscosity smooths it.  Pressure
keeps the motion incompressible and couples every local change to the
surrounding field.  The regularity problem is decided by whether nonlinear
sharpening can outrun viscous smoothing before some finite time \(T_*\).

The danger appears when the surrounding flow stretches a vortex.  Follow a
material portion of the tube as it lengthens: incompressibility preserves its
volume, so its cross-section contracts.  The core becomes thinner.  The axial
strain also acts on the rotation inside that core.  Write
\[
\omega:=\nabla\times u,
\qquad
S:=\frac12\bigl(\nabla u+(\nabla u)^{\mathsf T}\bigr).
\]
The local angular velocity is \(\omega/2\); \(S\) describes the stretching
and contraction.  Taking the curl of Navier--Stokes gives
\[
(\partial_t+u\!\cdot\nabla)\omega
=S\omega+\nu\Delta\omega.
\]
When the vorticity points along an extensional direction of \(S\), the term
\(S\omega\) points with it and strengthens the rotation.  Viscosity is acting
throughout, diffusing vorticity through \(\nu\Delta\omega\).  In the sharpening
scenario, stretching gains on diffusion and the fluid spins faster as the
core narrows.

That faster spin carries sharper rotational velocity gradients.  The
antisymmetric part of the velocity gradient is
\[
A:=\frac12\bigl(\nabla u-(\nabla u)^{\mathsf T}\bigr),
\qquad
|A|^2=\frac12|\omega|^2\leq|\nabla u|^2.
\]
Over the whole fluid, multiplying the vorticity equation by \(\omega\) and
integrating gives
\[
\frac12\frac{d}{dt}\|\omega(t)\|_2^2
=\int_{\mathbb R^3}\omega\!\cdot S\omega\,dx
-\nu\|\nabla\omega(t)\|_2^2.
\]

Nothing here requires the kinetic energy itself to become infinite.  In fact,
\[
\frac12\frac{d}{dt}\|u(t)\|_2^2
+\nu\|\nabla u(t)\|_2^2=0.
\]
The total energy decreases.  The unresolved possibility is concentration:
finite energy can be carried by thinner regions while the velocity gradients
inside those regions grow.  A vortex can be losing energy and still be
sharpening near its core.

So we shrink the picture and ask whether the balance changes.  For
\(\lambda>0\), set
\[
u_\lambda(x,t)=\lambda u(\lambda x,\lambda^2t),
\qquad
p_\lambda(x,t)=\lambda^2p(\lambda x,\lambda^2t).
\]
Every part of the equation picks up one common factor:
\[
\begin{aligned}
\partial_tu_\lambda
  &=\lambda^3(\partial_tu)(\lambda x,\lambda^2t),
&
(u_\lambda\!\cdot\nabla)u_\lambda
  &=\lambda^3((u\!\cdot\nabla)u)(\lambda x,\lambda^2t),\\
\nabla p_\lambda
  &=\lambda^3(\nabla p)(\lambda x,\lambda^2t),
&
\nu\Delta u_\lambda
  &=\lambda^3\nu(\Delta u)(\lambda x,\lambda^2t).
\end{aligned}
\]
Viscosity does not acquire an extra power of \(\lambda\) that makes it win
automatically in a smaller vortex.  Nonlinear transport, pressure, local
acceleration, and diffusion remain matched by scaling.

The energy balance can miss precisely this concentration.  Under the rescaling,
velocity squared contributes \(\lambda^2\), while spatial volume contributes
\(\lambda^{-3}\), giving
\[
\|u_\lambda(\cdot,t)\|_2^2
=\lambda^{-1}\|u(\cdot,\lambda^2t)\|_2^2.
\]
The concentrating copy becomes smaller in energy while its spatial changes
become steeper.  A spatial derivative contributes another factor
\(\lambda\), so measuring derivatives changes our sensitivity to the narrowing
core.  We want the height at which those effects cancel.  The Sobolev norm
\(\|\,\cdot\,\|_{\dot H^s}\) measures \(s\) spatial derivatives in \(L^2\),
including fractional orders.  Its scaling is
\[
\|u_\lambda(\cdot,t)\|_{\dot H^s}
=\lambda^{s-1/2}\|u(\cdot,\lambda^2t)\|_{\dot H^s},
\]
concentration lowers the norm below \(s=\tfrac12\), raises it above
\(s=\tfrac12\), and leaves it unchanged at
\[
s=\frac12.
\]
This is the critical height.  It is not a claim that nonlinear production and
viscous dissipation are numerically equal at each moment.  It says that scale
alone gives neither one an advantage there.

Now keep following the vortex as its core narrows.  If its radius is
\(\ell\), a second spatial derivative across that core has size
\(\ell^{-2}\) relative to the velocity change it sees.  The viscous term acts
on the corresponding interval
\[
|\nu\Delta u|
\sim\frac{\nu}{\ell^2}|u|,
\qquad
\tau_\ell\sim\frac{\ell^2}{\nu}.
\]
If the core kept narrowing through
\(\ell_j=\rho^j\ell_0\) with \(0<\rho<1\), those intervals would satisfy
\[
\sum_{j=0}^{\infty}\frac{\ell_j^2}{\nu}
=\frac{\ell_0^2}{\nu}\sum_{j=0}^{\infty}\rho^{2j}
=\frac{\ell_0^2}{\nu(1-\rho^2)}<\infty.
\]
That is the Zeno picture inside the vortex: each thinner core asks the equation
to respond faster, and the resulting intervals can fit below one finite
terminal time.  It does not assert that a solution follows this cascade.  It
shows the failure pattern the proof has to prevent.

To turn this possible compression into a contradiction, suppose that the smooth
solution has a finite maximal lifespan \(T_*\).  The scale calculation also
points to another critical way to measure velocity in three dimensions:
\[
\|u_\lambda(\cdot,t)\|_{L^3}
=\|u(\cdot,\lambda^2t)\|_{L^3}.
\]
The \(L^3\) norm, like \(\dot H^{1/2}\), does not change under concentration.
The endpoint continuation theorem now gives the decisive fork.  If
\[
\sup_{0<t<T_*}\|u(t)\|_{L^3}<\infty,
\]
then the solution continues smoothly through \(T_*\).  That is impossible if
\(T_*\) is maximal, so the contrapositive requires
\[
\sup_{0<t<T_*}\|u(t)\|_{L^3}=\infty.
\]
Here velocity blow-up means escape of the global critical \(L^3\) norm; it
does not mean that every particle speed must diverge.

The first finite-time integral keeps the physical meaning visible:
\[
u(t)=u_0+\int_0^t\partial_su(s)\,ds.
\]
If the local acceleration \(\partial_tu\) stayed uniformly bounded in
\(L^3\), this integral would keep the velocity bounded on the finite interval.
So velocity escape forces the local acceleration to escape.  The next integral
is
\[
\partial_tu(t)
=\partial_tu(0)+\int_0^t\partial_s^2u(s)\,ds.
\]
If the jerk \(\partial_t^2u\) stayed bounded, acceleration would stay bounded
and the first integral would again keep velocity bounded.  This is the
velocity--acceleration--jerk insight in plain language.  In order for a
singularity to form, velocity has to blow up.  In order for
that to happen, its acceleration would also have to blow up, which means its
jerk has to blow up, and so on, and so on:
\[
u,\qquad \partial_tu,\qquad \partial_t^2u,\qquad\ldots
\]
The smooth datum gives finite initial responses at every order.  A uniform
\(L^3\) bound on any fixed higher response would descend through finitely
many such integrals to a bound for \(u\).  So every fixed temporal order must
escape in \(L^3\) on \((0,T_*)\).

These responses are not added from outside.  Acceleration is already written
into Navier--Stokes.  Substitute the equation for it and differentiate: jerk
appears.  Differentiate again and the next response appears.  The first three
rungs are
\[
\begin{aligned}
\partial_tu
&=\nu\Delta u-(u\!\cdot\nabla)u-\nabla p,\\[4pt]
\partial_t^2u
&=\nu\Delta\partial_tu
  -(\partial_tu\!\cdot\nabla)u
  -(u\!\cdot\nabla)\partial_tu
  -\nabla\partial_tp,\\[4pt]
\partial_t^3u
&=\nu\Delta\partial_t^2u
  -(\partial_t^2u\!\cdot\nabla)u
  -2(\partial_tu\!\cdot\nabla)\partial_tu\\
&\hspace{2.6em}
  -(u\!\cdot\nabla)\partial_t^2u
  -\nabla\partial_t^2p.
\end{aligned}
\]
The complete pattern is
\[
\partial_t^{r+1}u
=\nu\Delta\partial_t^ru
-\sum_{a=0}^{r}\binom ra
  (\partial_t^au\!\cdot\nabla)\partial_t^{r-a}u
-\nabla\partial_t^rp.
\]

Only now do we compress the notation:
\[
V_r:=\partial_t^ru,
\qquad
Q_r:=\partial_t^rp.
\]
Pressure has to travel with every response.  Taking the divergence of
Navier--Stokes and using \(\nabla\!\cdot u=0\) gives
\[
-\Delta p=\partial_i\partial_j(u_i u_j).
\]
Define the Riesz transform by
\[
R_i:=\partial_i(-\Delta)^{-1/2}.
\]
With the decaying pressure normalization and repeated spatial indices summed,
this convention gives
\(p=R_iR_j(u_i u_j)\).  Differentiating that product \(r\) times recovers
the pressure response from the velocity responses.  The compact tower is
\[
\begin{aligned}
Q_r
&=R_iR_j\sum_{a=0}^{r}\binom ra V_{a,i}V_{r-a,j},\\
V_{r+1}
&=\nu\Delta V_r
  -\sum_{a=0}^{r}\binom ra(V_a\!\cdot\nabla)V_{r-a}
  -\nabla Q_r.
\end{aligned}
\]

Now the connection to spatial smoothness is visible.  A hypothetical
singularity runs forward through rougher and faster temporal responses.  The
proof reads the viscous relation backward.  Once a finite response and its
transport--pressure terms are controlled, the equation isolates
\(\Delta V_r\) and returns two spatial derivatives to the preceding response.
Repeating a finite number of those controlled steps would raise the Sobolev
height available to the original velocity.  Higher temporal order points toward
higher spatial order, but making that reversal useful requires bounds for the
full signed interaction carried at each step.

That brings us back to the contest that opened the section.  At critical
height, the energy pairing separates nonlinear production from viscous removal
without deleting either one.  Let
\[
\Lambda:=(-\Delta)^{1/2},
\qquad
E_s(t):=\frac12\|\Lambda^su(t)\|_2^2,
\qquad
H(t):=E_{1/2}(t).
\]
Pairing Navier--Stokes at height \(s\) gives
\[
E_s'=\mathcal P_s-2\nu E_{s+1},
\]
where the complete signed production is
\[
\mathcal P_s
:=-\bigl\langle
\Lambda^su,\,
\Lambda^s\bigl((u\!\cdot\nabla)u+\nabla p\bigr)
\bigr\rangle.
\]
Repeated differentiation at \(s=\tfrac12\) shows how viscous removal exposes
the next spatial energy while the signed production differentiates alongside
it:
\[
\begin{aligned}
H'
&=\mathcal P_{1/2}-2\nu E_{3/2},\\
H''
&=\partial_t\mathcal P_{1/2}
  -2\nu\mathcal P_{3/2}
  +(2\nu)^2E_{5/2},\\
H'''
&=\partial_t^2\mathcal P_{1/2}
  -2\nu\partial_t\mathcal P_{3/2}
  +(2\nu)^2\mathcal P_{5/2}
  -(2\nu)^3E_{7/2}.
\end{aligned}
\]
At every finite order,
\[
H^{(n)}
=(-2\nu)^nE_{n+1/2}
+\sum_{j=0}^{n-1}(-2\nu)^j
 \partial_t^{\,n-1-j}\mathcal P_{j+1/2}.
\]
The \(n\)-th temporal derivative has exposed the spatial energy at height
\(n+\tfrac12\), together with the production terms that must be controlled.

Now the reversal has its payoff.  Together with the \(L^2\) energy already
controlled, a bound at a higher spatial energy gives the corresponding
\(H^s\) bound.  For any requested classical derivative order \(k\), it would
be enough to reach
\[
s>k+\frac32,
\qquad
H^s(\mathbb R^3)\hookrightarrow C_b^k(\mathbb R^3).
\]
If every finite spatial height can be controlled, then no finite derivative
order is left out:
\[
\bigcap_{m<\infty}H^m(\mathbb R^3)
=H^\infty(\mathbb R^3)
\hookrightarrow C^\infty(\mathbb R^3).
\]
The temporal cascade that seemed to demand worsening behavior has exposed a
route to complete spatial smoothness.  This is the insight; its condition is
still to be proved.  The differentiated identities have exposed coordinates,
not made them finite by rearrangement.  Fix a finite
temporal depth \(N\) and spatial height \(M\).  One step in the recurrence
advances the temporal response by one and passes through two spatial derivatives
in \(\Delta V_r\).  That is why the finite estimate pairs height \(M+1\)
of \(V_r\) with height \(M-1\) of \(V_{r+1}\).

The whole-terminal rectangle theorem proves the required bound:
\[
\begin{aligned}
\mathfrak R_{N,M}(T)
&:=\sum_{r=0}^{N}
\|V_r\|_{L^2(0,T;H^{M+1})}^{2}+
\sum_{r=0}^{N}
\|V_{r+1}\|_{L^2(0,T;H^{M-1})}^{2},\\
\sup_{0<T<T_*}\mathfrak R_{N,M}(T)&<\infty.
\end{aligned}
\]
Here \(N\) and \(M\) are fixed before \(T\) approaches \(T_*\).  This is
the bound that survives every cutoff; smoothness on each shorter interval
would not supply it.  Let \(T\uparrow T_*\).  Each term is an integral of a
nonnegative squared norm, so monotone convergence gives, for
\(0\le r\le N\),
\[
V_r\in L^2(0,T_*;H^{M+1}),
\qquad
V_{r+1}\in L^2(0,T_*;H^{M-1}),
\]
and the pressure formula gives the finite companion estimate
\[
\|Q_r\|_{L^1(0,T_*;H^M)}
\leq C_{M,r}\sum_{a=0}^{r}
\|V_a\|_{L^2(0,T_*;H^{M+1})}
\|V_{r-a}\|_{L^2(0,T_*;H^{M+1})}<\infty.
\]
This uses the Riesz-transform bound, finite-order Sobolev products, and
Cauchy--Schwarz in time.  Pressure remains inside the finite record.
The velocity estimate controls \(V_r\) together with its time derivative
\(V_{r+1}=\partial_tV_r\).  That second row restricts how quickly the first can
change.  Their spatial heights sit on either side of \(H^M\) in the Hilbert
triple
\[
H^{M+1}\subset H^M\subset H^{M-1}.
\]
For \(0\le a<b<T_*\), differentiate the squared \(H^M\) norm and pair the
two rows.  Cauchy--Schwarz in space and time gives
\[
\left|
\|V_r(b)\|_{H^M}^2-\|V_r(a)\|_{H^M}^2
\right|
\leq
2\|V_{r+1}\|_{L^2(a,b;H^{M-1})}
 \|V_r\|_{L^2(a,b;H^{M+1})}.
\]
Both factors on the right tend to zero as the time interval shrinks, because
their squared norms are integrable through \(T_*\).  A large rise in the
middle norm cannot be hidden in an arbitrarily short interval.  In particular,
the squared norm has a finite limit as time approaches \(T_*\).
This estimate controls the norm; the Lions--Magenes trace theorem uses the
full adjacent-row hypothesis to supply strong continuity of the field at the
middle height, including the endpoint:
\[
V_r\in C([0,T_*];H^M).
\]
This is one finite rectangle.  Its constant may depend on \(N\) and \(M\).
The proof never asks one constant to control infinitely many derivatives at
once.  It proves the required statement for every finite rectangle.

Those rectangles are not unrelated solutions.  Every coordinate belongs to
one velocity history generated by \(u_0\).  When two rectangles overlap, their
shared coordinates agree.  Compatibility lets us assemble the entire family:
\[
u_0
\longrightarrow
\{\mathcal R_{N,M}[u_0;T_*]\}_{N,M<\infty}
\longrightarrow
\mathcal I[u_0]
:=\bigcap_{r,m<\infty}H^r(0,T_*;H^m(\mathbb R^3)),
\qquad
u\in\mathcal I[u_0].
\]
No infinite-order estimate has been smuggled into this step.  Each coordinate
came from a finite theorem; the intersection records that all of them belong
to one compatible fluid.

Now take the zeroth temporal row.  For every finite \(m\), the rectangle pair
and the equation place
\[
u\in C([0,T_*];H^m_\sigma).
\]
All of these traces meet at one terminal value,
\[
u_*:=u(T_*)
\in\bigcap_{m<\infty}H^m_\sigma
=H^\infty_\sigma
\hookrightarrow C^\infty.
\]
At \(T_*\), the equation converts the completed spatial information back into
time information.  Starting from \(V_0^*=u_*\), the pressure formula gives
\(Q_0^*\), and the recurrence gives \(V_1^*\).  Repeating this at every finite
order produces the pressure-complete terminal jet.  Smooth local existence
restarts from \(u_*\); uniqueness joins that continuation to the original
history.  This extends the solution past \(T_*\), contradicting finite
maximality.  Hence
\[
T_*=\infty.
\]

After continuation, fix any finite time \(T\) and write \(I=(0,T)\).  The
construction on that interval gives
\[
u\in\mathcal I_T[u_0]
:=\bigcap_{a,m<\infty}H^a(I;H^m(\mathbb R^3)).
\]
The familiar estimates are read from this finite-interval intersection instead
of proved one by one.  Every time space below refers to \(I\), not to the whole
infinite future.  Given a finite response order \(r\), a spatial
derivative order \(k\), and \(2\le p\le\infty\), choose
\[
m>k+\frac32-\frac3p.
\]
That finite coordinate preserves the requested response order:
\[
\begin{aligned}
u\in\mathcal I_T[u_0]
&\Longrightarrow u\in H_t^{r+1}H_x^m
\Longrightarrow V_r\in H_t^1H_x^m\\
&\Longrightarrow V_r\in C_tH_x^m
\Longrightarrow V_r\in L_t^qW_x^{k,p}
\qquad(1\le q\le\infty).
\end{aligned}
\]
At \(r=0\), the endpoint embedding
\(H^1(\mathbb R^3)\hookrightarrow L^6(\mathbb R^3)\), followed by the strict
embeddings above with \(m=2\) and \(m=3\), recovers, among many others,
\[
u\in L_t^\infty L_x^6,
\qquad
u\in L_t^\infty L_x^\infty,
\qquad
u\in L_t^\infty W_x^{1,\infty}.
\]
These are not separate achievements.  Each is one finite readout of
\(\mathcal I_T[u_0]\).

The whole motion is
\[
\boxed{
u_0
\longrightarrow \{(V_r,Q_r)\}_{r\ge0}
\longrightarrow \{\mathcal R_{N,M}[u_0]\}_{N,M<\infty}
\longrightarrow \mathcal I[u_0]
\longrightarrow u_*\in H^\infty
\longrightarrow \text{restart}
\longrightarrow T_*=\infty.
}
\]

\endgroup
